% Options for packages loaded elsewhere
\PassOptionsToPackage{unicode}{hyperref}
\PassOptionsToPackage{hyphens}{url}
\documentclass[
  hebrew,
]{article}
\usepackage{xcolor}
\usepackage{amsmath,amssymb}
\setcounter{secnumdepth}{-\maxdimen} % remove section numbering
\usepackage{iftex}
\ifPDFTeX
  \usepackage[T1]{fontenc}
  \usepackage[utf8]{inputenc}
  \usepackage{textcomp} % provide euro and other symbols
\else % if luatex or xetex
  \usepackage{unicode-math} % this also loads fontspec
  \defaultfontfeatures{Scale=MatchLowercase}
  \defaultfontfeatures[\rmfamily]{Ligatures=TeX,Scale=1}
\fi
\usepackage{lmodern}
\ifPDFTeX\else
  % xetex/luatex font selection
  \setmainfont[]{DejaVu Sans}
\fi
% Use upquote if available, for straight quotes in verbatim environments
\IfFileExists{upquote.sty}{\usepackage{upquote}}{}
\IfFileExists{microtype.sty}{% use microtype if available
  \usepackage[]{microtype}
  \UseMicrotypeSet[protrusion]{basicmath} % disable protrusion for tt fonts
}{}
\makeatletter
\@ifundefined{KOMAClassName}{% if non-KOMA class
  \IfFileExists{parskip.sty}{%
    \usepackage{parskip}
  }{% else
    \setlength{\parindent}{0pt}
    \setlength{\parskip}{6pt plus 2pt minus 1pt}}
}{% if KOMA class
  \KOMAoptions{parskip=half}}
\makeatother
\ifLuaTeX
\usepackage[bidi=basic,shorthands=off,provide=*]{babel}
\else
\usepackage[bidi=default,shorthands=off,provide=*]{babel}
\fi
\ifPDFTeX
\else
\babelfont{rm}[]{DejaVu Sans}
\fi
\ifLuaTeX
  \usepackage{selnolig} % disable illegal ligatures
\fi
\setlength{\emergencystretch}{3em} % prevent overfull lines
\providecommand{\tightlist}{%
  \setlength{\itemsep}{0pt}\setlength{\parskip}{0pt}}
\ifPDFTeX
  \TeXXeTstate=1
  \newcommand{\RL}[1]{\beginR #1\endR}
  \newcommand{\LR}[1]{\beginL #1\endL}
  \newenvironment{RTL}{\beginR}{\endR}
  \newenvironment{LTR}{\beginL}{\endL}
\fi
\ifluatex
  \newcommand{\RL}[1]{\bgroup\textdir TRT#1\egroup}
  \newcommand{\LR}[1]{\bgroup\textdir TLT#1\egroup}
  \newenvironment{RTL}{\textdir TRT\pardir TRT\bodydir TRT}{}
  \newenvironment{LTR}{\textdir TLT\pardir TLT\bodydir TLT}{}
\fi
\usepackage{bookmark}
\IfFileExists{xurl.sty}{\usepackage{xurl}}{} % add URL line breaks if available
\urlstyle{same}
% fallback for those not using the hyperref driver hyperxmp:
\makeatletter
\@ifundefined{xmpquote}{\newcommand{\xmpquote}[1]{#1}}{}
\makeatother
\hypersetup{
  pdflang={he},
  hidelinks,
  pdfcreator={LaTeX via pandoc}}

\author{}
\date{}

\begin{document}

\section{מרק ולצר, קורות
חיים}\label{ux5deux5e8ux5e7-ux5d5ux5dcux5e6ux5e8-ux5e7ux5d5ux5e8ux5d5ux5ea-ux5d7ux5d9ux5d9ux5dd}

\subsection{פרטי
התקשרות}\label{ux5e4ux5e8ux5d8ux5d9-ux5d4ux5eaux5e7ux5e9ux5e8ux5d5ux5ea}

\begin{itemize}
\tightlist
\item
  972-50-5665636+ (נייד)
\item
  mark.veltzer@gmail.com
\item
  \url{https://github.com/veltzer}
\end{itemize}

\subsection{כישורים
עיקריים}\label{ux5dbux5d9ux5e9ux5d5ux5e8ux5d9ux5dd-ux5e2ux5d9ux5e7ux5e8ux5d9ux5d9ux5dd}

\begin{itemize}
\tightlist
\item
  מדריך בכיר (לימדתי יותר מ-150 קורסים שונים)
\item
  מדריך AWS
\item
  לינוקס משובץ (Embedded)
\item
  לינוקס בזמן אמת (Real Time)
\item
  תכנות בלינוקס
\item
  מומחה יוניקס
\item
  C/C++/Java/Perl/Python/Php/Scala/JavaScript/Bash
\item
  ג'נרליסט
\item
  למידת מכונה
\item
  מדעי הנתונים
\item
  אלגוריתמים
\item
  אופטימיזציה
\item
  אלגוריתמים גנטיים
\end{itemize}

\subsection{שפות}\label{ux5e9ux5e4ux5d5ux5ea}

\begin{itemize}
\tightlist
\item
  אנגלית (שפת אם או דו-לשוני)
\item
  עברית (שפת אם או דו-לשוני)
\item
  רוסית (רמה גבוהה)
\item
  ספרדית (מתחיל)
\end{itemize}

\subsection{קורסים
והסמכות}\label{ux5e7ux5d5ux5e8ux5e1ux5d9ux5dd-ux5d5ux5d4ux5e1ux5deux5dbux5d5ux5ea}

\begin{itemize}
\tightlist
\item
  Games without Chance: Combinatorial Game Theory, תומאס מורלי, Georgia
  Tech, דרך Coursera
\item
  קיצור תולדות האנושות, יובל נח הררי, האוניברסיטה העברית בירושלים, דרך
  Coursera
\item
  אלתור ג'אז, גארי ברטון, Berklee College of Music, דרך Coursera
\item
  מדריך מוסמך של Sun Microsystems ל-Solaris ול-Java מתקדם, כיום Oracle
\item
  מדריך מוסמך ל-MySQL, כיום Oracle
\item
  AWS: ‏Associate Architect
\item
  AWS: מדריך מוסמך: Architecting
\item
  AWS: ‏Associate Advanced Architect
\item
  AWS: מדריך מוסמך: Architecting, רמה מקצועית (Professional)
\item
  AWS: ‏Devops on AWS
\end{itemize}

\subsection{השכלה}\label{ux5d4ux5e9ux5dbux5dcux5d4}

\begin{itemize}
\tightlist
\item
  אוניברסיטת תל אביב, BSc, מתמטיקה (1990 - 1993)
\item
  בית הספר רימון לג'אז, ביצוע ג'אז (2004 - 2009)
\item
  האוניברסיטה הפתוחה, BA, פילוסופיה והיסטוריה (2010 - 2012)
\end{itemize}

\subsection{פרסומים}\label{ux5e4ux5e8ux5e1ux5d5ux5deux5d9ux5dd}

\begin{itemize}
\item
  Democracy Unleashed, A 21st Century Manifesto (עם דיוויד אלכסנדר)

  ספר על עידן הדמוקרטיה הישירה המתקרב, המציג סקירה של שיטות הממשל שקדמו
  לו.
\end{itemize}

\subsection{תקציר}\label{ux5eaux5e7ux5e6ux5d9ux5e8}

מתכנת hands-on בשפות רבות. אני עובד בלפחות 5-6 שפות ולימדתי בערך 12. אני
כותב קוד החל מ-C ברמת הקרנל למימוש דרייברים ועד Scala ברמת Enterprise
ורוב השפות שביניהן. אני גם מדריך טוב והעברתי מאות קורסים לחברות
טכנולוגיה רבות בישראל וברחבי העולם. אני ארכיטקט ולא רק מתכנת, והייתי
מעורב בתכנון הארכיטקטורה של כמה מוצרים: מערכת בית חכם עבור אינטל, אתר
בעל יכולת סקיילינג גבוהה מאוד עבור יאהו, אפליקציית זמן אמת עבור רפאל,
אפליקציית עיבוד סקיילבילית על AWS עבור לקוחות פרטיים שלי ועוד. הדרכתי
בערך 150 קורסים שונים, מבחינת סילבוס, כל אחד מהם כמה פעמים, במהלך העשור
האחרון, בדרך כלל לאנשי מקצוע בתעשייה. מקורסים מתודולוגיים כמו תכנות
אג'ילי ועד קורסים טכניים מאוד, כמו תכנות מרובה ליבות על מעבדי אינטל, דרך
כלים כמו git ועד קורסים פילוסופיים כמו הכנסת קוד פתוח לחברות. אני לומד
עצמאית את רוב הנושאים שאני מדריך ומשתמש בהם. יש לי רקע מתמטי חזק (לא
השלמתי תואר שני רק מתוך עצלות). ניהלתי צוותים בפרויקטים שונים אבל אני
יכול גם לעבוד לבד. כשאני מנהל צוות אני משתדל להימנע מפרקטיקות בירוקרטיות
ולהיות חלק מהצוות. תרמתי לפרויקטים רבים של תוכנה חופשית וקוד פתוח כמו
קרנל הלינוקס, וכיום אני המתחזק של GNU grep(1). אני מתכנת גם כתחביב
בפרויקטים אישיים שאפשר למצוא ב-github. תחומי העניין האהובים עליי הם קרנל
הלינוקס, ארכיטקטורת UNIX ובניית מערכות לינוקס משובצות. אני גם עושה את
ניהול המערכת שלי בעצמי ולימדתי קורסים רבים בניהול מערכות. אני מוסמך
ב-mysql, solaris, לינוקס, AWS ועוד. אני גם DBA מקצועי של MySQL, הצעתי
תמיכה מקצועית בתחום ולימדתי כמה קורסי MySQL.

הייתי ה-associate instructor הראשון של AWS בישראל.

לאחרונה אני עוסק במדעי הנתונים, למידת מכונה ו-NLP על נתונים בגודל בינוני
(כמה טרה-בייטים).

\subsection{פעילות
קהילתית}\label{ux5e4ux5e2ux5d9ux5dcux5d5ux5ea-ux5e7ux5d4ux5d9ux5dcux5eaux5d9ux5ea}

\begin{itemize}
\tightlist
\item
  חבר GNU משנת 2001 (veltzer@gnu.org)
\item
  הייתי המתחזק של GNU grep(1) מ-2013 עד 2020
\item
  מודולי python רבים (ראו \url{https://pypi.org/user/veltzer/})
\item
  מאגרים רבים ב-github (ראו \url{https://github.com/veltzer})
\end{itemize}

\subsection{ניסיון
תעסוקתי}\label{ux5e0ux5d9ux5e1ux5d9ux5d5ux5df-ux5eaux5e2ux5e1ux5d5ux5e7ux5eaux5d9}

\begin{itemize}
\tightlist
\item
  יועץ עצמאי

  \begin{itemize}
  \tightlist
  \item
    יועץ בכיר וכותב קוד קשוח
  \item
    ינואר 2021 - היום
  \item
    חולון, ישראל
  \item
    אני מייעץ לחברות גדולות מאוד רבות, בעיקר בנושא לינוקס משובץ. רפאל,
    התעשייה האווירית, אינטל, אתם מבינים את התמונה.
  \end{itemize}
\item
  LiteBC

  \begin{itemize}
  \tightlist
  \item
    ארכיטקט בכיר ומתכנת
  \item
    אוקטובר 2022 - 2024 (שנתיים)
  \item
    יקנעם, ישראל
  \item
    עשיתי הרבה עבודת תשתיות והכוונה לגבי אופן פיתוח המחקר והאסטרטגיה של
    החברה.
  \end{itemize}
\item
  Sqlink

  \begin{itemize}
  \tightlist
  \item
    מדריך בכיר
  \item
    אוגוסט 2019 - היום
  \item
    תל אביב, ישראל
  \item
    אני מלמד קורסים שונים ברמות מומחיות שונות. העברתי גם בוטקאמפ ארוך
    למתכנתי לינוקס משובץ מתחילים.
  \end{itemize}
\item
  AgileSparks

  \begin{itemize}
  \tightlist
  \item
    יועץ טכני בכיר
  \item
    אוקטובר 2019 - היום
  \item
    תל אביב, ישראל
  \item
    אני מייעץ לחברות שונות בנוגע לשימוש נכון בטכנולוגיה ולארכיטקטורה
    בנושאים רבים, ובפרט שימוש יעיל במשאבי ענן ציבורי כגון AWS.
  \end{itemize}
\item
  ג'ון ברייס

  \begin{itemize}
  \tightlist
  \item
    מדריך בכיר
  \item
    אוגוסט 2013 - היום
  \item
    תל אביב, ישראל
  \item
    אני מלמד כל מה שקשור ללינוקס, java, mysql, תכנות מערכות unix, קרנל
    לינוקס, Python, Perl, PHP, Javascript, שרת האינטרנט Apache, ACE,
    JEE, Spring, Hibernate, Swing, Ruby ועוד קורסים רבים רבים אחרים.
  \end{itemize}
\item
  Interbit T\&C Ltd.

  \begin{itemize}
  \tightlist
  \item
    מדריך בכיר
  \item
    2006 - היום
  \item
    תל אביב, ישראל
  \item
    אני מלמד כל מה שקשור ל-java, mysql, תכנות מערכות unix, קרנל לינוקס,
    Python, Perl, PHP, Javascript, שרת האינטרנט Apache, ACE, JEE,
    Spring, Hibernate, Swing ועוד קורסים רבים רבים אחרים.
  \end{itemize}
\item
  Twiggle

  \begin{itemize}
  \tightlist
  \item
    מהנדס תוכנה בכיר
  \item
    אוגוסט 2016 - ינואר 2020 (3 שנים ו-6 חודשים)
  \item
    תל אביב, ישראל
  \item
    סטארטאפ בתחום הבינה המלאכותית. עבדתי בצוות ה-CTO על פרויקטים שונים
    בחזית הטכנולוגיה. באופן כללי עסקתי ב-big data, בינה מלאכותית, NLP,
    למידת מכונה ואלגוריתמים.
  \end{itemize}
\item
  רפאל

  \begin{itemize}
  \tightlist
  \item
    מפתח ליבה
  \item
    2006 - 2011 (5 שנים)
  \item
    הצפון, ישראל
  \item
    עשיתי הנדסת ליבה בפרויקט מסווג. פנימי קרנל לינוקס, C++ ועוד.
  \end{itemize}
\item
  Hinbit Development

  \begin{itemize}
  \tightlist
  \item
    CTO
  \item
    2005 - אוגוסט 2016 (11 שנים)
  \item
    רעננה, ישראל
  \item
    עשינו עבודות קבלנות משנה בתחומים שונים עם התמחות במיגרציה ללינוקס
    מכל דבר. שימו לב - הכוונה למיגרציה אל לינוקס, לא מ- לינוקס. אנחנו
    כותבים דרייברים ללינוקס, עוזרים בסוגיות זמן אמת, מציעים ייעוץ,
    קורסים, שירותי מיגרציית קוד מפלטפורמות אחרות ועוד. אנחנו עושים גם
    פרויקטי web מסוגים שונים. עבדנו עם כמה מהחברות הגדולות בישראל:
    אמדוקס, אינטל, IBM, צ'קפוינט.
  \end{itemize}
\item
  Qlusters

  \begin{itemize}
  \tightlist
  \item
    מהנדס קרנל ו-userspace
  \item
    2003 - 2005 (שנתיים)
  \item
    תל אביב, ישראל
  \item
    Qlusters עשתה עבודה בתוך הקרנל לתמיכה ב-Mosix (קרנל אחד, מכונות
    רבות) עבור לקוחות עסקיים. עבדתי גם בקרנל וגם ב-user space.
  \end{itemize}
\item
  Optivera

  \begin{itemize}
  \tightlist
  \item
    ג'נרליסט
  \item
    2001 - 2003 (שנתיים)
  \item
    תל אביב, ישראל
  \item
    זה היה סטארטאפ בתחום הבינה המלאכותית שניסה להשתמש באלגוריתמים גנטיים
    לאופטימיזציה של ניתוב רשתות.
  \end{itemize}
\item
  Schema

  \begin{itemize}
  \tightlist
  \item
    ג'נרליסט
  \item
    1997 - 2001 (4 שנים)
  \item
    הרצליה, ישראל
  \item
    Schema הייתה חנות פתוחה לאלגוריתמים גנטיים כשהגעתי אליה. עבדנו על
    לינוקס עם C++ לאופטימיזציה של בעיות NP שלמות שעולות בתעשיות שונות.
    עבדתי על בעיות שונות: העמסת אוניות משא, איזון טילים, הקצאת תדרים
    סלולריים ועוד. בסופו של דבר החברה עברה לעסוק בסלולר בלבד ואני הייתי
    מהמפתחים הבכירים של המוצר המרכזי של החברה.
  \end{itemize}
\item
  צה''ל, משרד הביטחון

  \begin{itemize}
  \tightlist
  \item
    פיתוח C++ וחקר ביצועים
  \item
    1993 - 1997 (5 שנים)
  \item
    איפשהו
  \item
    עבדתי בזרוע חקר הביצועים של משרד הביטחון. העבודה כללה C++, יוניקס,
    Solaris, לוגיקה עמומה, סימולציות, סטטיסטיקה, מיפוי, עיצוב ממשקי
    משתמש ועוד דברים שאיני יכול לדבר עליהם.
  \end{itemize}
\end{itemize}

\end{document}
